\documentclass{article}
\usepackage{graphicx} % Required for inserting images

\title{A Ciência do Sagrado e a Estética da Exatidão: Transdisciplinaridade, Observação da Natureza e Expressão Artística na Sociedade Maia}
\author{DAVID e CARLOS}
\date{November 2026}

\begin{document}

\maketitle


\section{Introdução}

A maneira como compreendemos e organizamos o conhecimento está profundamente relacionada à sociedade na qual estamos inseridos. Na contemporaneidade, estamos habituados a reconhecer a ciência, a arte, a religião e a tecnologia como campos distintos da atividade humana, cada qual com suas próprias linguagens, métodos e formas de expressão. Embora suas relações e interseções sejam frequentemente reconhecidas, essa divisão permanece bastante presente na maneira como interpretamos o mundo, produzimos conhecimento e estudamos a história das sociedades. Entretanto, ao observarmos outras civilizações e contextos históricos, torna-se necessário questionar até que ponto essas classificações correspondem a formas universais de organização da experiência humana ou se são, elas próprias, construções históricas e culturais.

Essa questão adquire especial relevância quando nos voltamos às sociedades pré-colombianas das Américas. Frequentemente, os conhecimentos produzidos por esses povos chegam até nós sob narrativas que oscilam entre duas perspectivas igualmente problemáticas: de um lado, sua desvalorização diante dos parâmetros estabelecidos pela tradição científica ocidental; de outro, sua idealização como realizações extraordinárias, frequentemente associadas a conhecimentos misteriosos ou a uma suposta precisão quase inexplicável. Ambas as perspectivas podem dificultar a compreensão de como esses conhecimentos foram efetivamente produzidos, transmitidos e incorporados à organização da vida social. Como argumenta Johanna Broda (2012), compreender a produção de conhecimento nas sociedades pré-hispânicas exige situá-la em seu contexto histórico e cultural, considerando as relações entre observação da natureza, cosmologia, religião e organização social, sem reduzi-la aos padrões institucionais da ciência moderna ou recorrer a uma idealização acrítica dessas civilizações.

Nesse sentido, estudar as sociedades antigas do continente americano representa não apenas uma oportunidade de reconhecer tradições intelectuais frequentemente marginalizadas nas narrativas convencionais da história do conhecimento, mas também de refletir criticamente sobre nossas próprias categorias de interpretação. Entre essas sociedades, a civilização maia oferece um campo particularmente rico para essa investigação. Seus conhecimentos matemáticos e calendáricos, a observação sistemática dos fenômenos celestes, os registros escritos, a arquitetura monumental e as manifestações artísticas constituem diferentes evidências de uma produção intelectual complexa, profundamente relacionada à cosmologia, às práticas religiosas e à organização política. Mais do que identificar realizações que possam ser classificadas isoladamente como científicas ou artísticas segundo nossos critérios atuais, interessa compreender como essas diferentes dimensões participavam da construção e da transmissão de conhecimentos em seus contextos originais.

A partir dessa perspectiva, este ensaio busca investigar como arte e ciência, entendidas aqui como categorias analíticas contemporâneas, se relacionavam na sociedade maia, considerando os fatores culturais, sociais, políticos e cosmológicos que permeavam suas formas de produção de conhecimento. Parte-se da hipótese de que determinadas manifestações materiais maias não podem ser compreendidas adequadamente quando separadas em categorias rígidas, uma vez que podiam articular conhecimentos astronômicos, matemáticos e calendáricos a expressões visuais, práticas rituais e mecanismos de organização social. Isso não significa, contudo, pressupor que os maias reconhecessem essas atividades como disciplinas integradas nos mesmos termos em que as concebemos atualmente, nem que toda produção material possuísse necessariamente uma função científica, artística ou religiosa simultânea.

Para desenvolver essa investigação, são examinados três conjuntos de evidências: o Castillo de Kukulcán, em Chichén Itzá, cuja arquitetura e relação com os fenômenos solares permitem discutir possíveis articulações entre observação astronômica, expressão visual e práticas rituais; as tabelas astronômicas do Códice de Dresden, que evidenciam a sistematização de conhecimentos calendáricos e a organização visual de registros numéricos; e os murais de Bonampak, por meio dos quais se investiga a dimensão narrativa, simbólica e política das manifestações visuais maias. A análise procura distinguir evidências materiais documentadas, interpretações propostas pelos pesquisadores e hipóteses ainda debatidas, evitando atribuir aos objetos significados ou funções que não possam ser sustentados pelas fontes disponíveis. Por fim, busca-se compreender o que esses casos revelam não apenas sobre a produção de conhecimento na sociedade maia, mas também sobre os limites das classificações modernas que utilizamos para interpretar as diferentes manifestações da experiência humana.

\section{Marco teórico: observação da natureza, instituição e contexto cultural}

Johanna Broda propõe que a observação da natureza seja tratada como atividade histórica, não como uma coleção de curiosidades astronômicas. Sua definição é operacional: “la observación sistemática y repetida a través del tiempo de los fenómenos naturales del medio ambiente que permite hacer predicciones y orientar el comportamiento social” (Broda, 2012, p. 107). A formulação reúne repetição, conhecimento acumulado, previsão e uso social. Ela permite reconhecer práticas empíricas sem presumir uma ciência moderna separada de religião, governo ou produção.

Essa perspectiva combate dois erros simétricos. O primeiro é o eurocentrismo que só chama de ciência uma prática institucionalizada segundo moldes modernos. O segundo é idealizar qualquer conhecimento antigo como ciência avançada ou sabedoria ecológica perfeita. Broda adverte contra “una idealización a-crítica de los logros de las culturas prehispánicas” (2012, p. 104) e assinala que conceitos como ciência e natureza precisam ser entendidos em sua concretização histórica (p. 106). Portanto, “ciência maia” será usada aqui como categoria analítica para conhecimentos empíricos e sistemáticos documentados, não como nome de uma instituição equivalente a uma universidade ou academia moderna.

O conceito de polivalência funcional ajuda a pensar instituições e espaços que podem reunir funções políticas, rituais, técnicas e comunicativas. Broda retoma a ideia de “fusión institucional” para explicar que, em sociedades antigas, esferas econômicas, sociais, políticas e ideológicas não devem ser estudadas como domínios isolados (2012, p. 106). A consequência metodológica não é declarar todo templo um observatório ou todo mural um calendário. É investigar se evidências diferentes mostram que uma mesma estrutura ou prática servia, em certos contextos, a finalidades sobrepostas.

Aveni formula uma cautela convergente para a arqueoastronomia: o registro arqueológico permite relacionar alinhamentos, calendário, ritual e poder, mas a função prática e a performativa não devem ser isoladas. Em muitos casos, lugares sagrados eram “stages upon which seasonally timed rituals were carried out” (Aveni, 2003, p. 175). Šprajc e Sánchez Nava concluem, de forma semelhante, que em Chichén Itzá “la función utilitaria de la astronomía estaba envuelta en el ritual e íntimamente relacionada con la vida social, la religión y la ideología política” (2013, p. 55). A convergência sustenta uma hipótese de articulação, mas não prova automaticamente a intenção de cada edifício.

Essa distinção é especialmente importante ao mobilizar Broda. Seu capítulo trata do México pré-hispânico e de Mesoamérica em escala ampla, não exclusivamente das cidades maias. A autora situa calendário e astronomia em processos regionais e instituições específicas; o ensaio usa seu método para formular perguntas, não para transferir diretamente conclusões sobre os mexicas a Chichén Itzá, Bonampak ou ao Dresden. A validade da ponte depende das pesquisas maias que documentam os objetos em questão.

\section{Estudos de caso}

\subsection{ Castillo de Kukulcán: arquitetura, luz e hipótese calendárica}

Chichén Itzá foi um centro político importante no norte das Terras Baixas durante o Clássico Terminal e o Pós-Clássico Inicial (Šprajc; Sánchez Nava, 2013, p. 33). A cronologia do sítio e a relação entre seus estilos arquitetônicos continuam debatidas (Milbrath, 1988, pp. 57–59). Por isso, o Castillo deve ser situado nessa sequência histórica complexa, sem ser tratado como obra de um único momento ou como expressão uniforme de toda a sociedade maia.

Na escadaria norte, a incidência solar produz, em determinadas épocas do ano, triângulos de luz e sombra sobre a balaustrada. Aveni descreve sete formas luminosas que lembram o padrão de uma cascavel e registra a associação visual com as cabeças de serpente esculpidas na base da escadaria (2003, p. 163). Essa descrição sustenta a existência de um efeito óptico e de elementos arquitetônicos que favorecem a leitura da serpente. Não demonstra, por si só, que os construtores tenham concebido o fenômeno como marcador equinocial ou como uma hierofania planejada; essas afirmações dizem respeito à intenção e exigem inferência histórica.

A cautela é reforçada pela reavaliação de Šprajc e Sánchez Nava. A amostra de 271 orientações em 87 sítios das Terras Baixas Maias fornece um contexto comparativo para estudar a arquitetura, em vez de tomar o Castillo isoladamente (2013, p. 35). Os autores observam que o efeito luminoso pode ser percebido durante vários dias em torno do equinócio e não permite identificar, apenas pela observação, uma data exata pretendida pelos construtores (p. 45). Portanto, a coincidência visual com a época equinocial não basta para provar que o edifício foi desenhado para assinalar o equinócio.

O estudo identifica padrões de orientações solares associados a datas separadas por intervalos calendáricos e propõe que observações do horizonte poderiam contribuir para acompanhar o ano sazonal e organizar atividades agrícolas e rituais (Šprajc; Sánchez Nava, 2013, pp. 31–32, 55). Essa interpretação regional não determina automaticamente a função do Castillo. Ela sustenta a possibilidade de que orientação, calendário e práticas sociais se relacionassem, mas a finalidade particular do edifício permanece uma questão interpretativa.

Milbrath propõe outras associações astronômicas para edifícios e imagens de Chichén Itzá, incluindo referências ao zênite solar, às Plêiades e a Vênus (1988, pp. 64–68). Sua leitura de oito cenas do Templo Superior dos Jaguares como possível referência ao ciclo de Vênus é uma hipótese iconográfica, não uma identificação inequívoca (pp. 66–67). A própria discussão registra problemas em interpretações anteriores que combinavam ciclos sem corresponder a uma configuração astronômica real (p. 67). Essas propostas ampliam as perguntas possíveis, mas não devem ser apresentadas como consenso nem como confirmação independente do fenômeno luminoso do Castillo.

O caso permite, assim, uma conclusão proporcional às fontes: há um efeito de luz e sombra observável e uma arquitetura que o torna visualmente legível; a relação com calendário, ritual e audiência é plausível no contexto arqueoastronômico mais amplo; a data equinocial precisa e a intenção específica não estão demonstradas. A articulação entre forma construída e fenômeno celeste é relevante para a análise estética do ensaio, desde que essa análise seja apresentada como leitura contemporânea, e não como prova de uma teoria maia da beleza.

\subsection{ Códice de Dresden: cálculo, imagem e transmissão}

O Códice de Dresden é um dos poucos manuscritos maias pré-hispânicos sobreviventes. Essa preservação restrita cria uma assimetria documental: o códice oferece evidência importante sobre determinados conhecimentos e formas de registro, mas não representa sozinho todas as comunidades, períodos ou práticas maias. Broda destaca a relevância dos códices de Dresden e Madrid para estudar conhecimentos calendáricos (2012, p. 111). Sua expressão “verdaderos tratados científicos” descreve o valor desses manuscritos para a pesquisa; não deve ser tomada como prova de que os próprios maias organizassem esse saber sob a categoria moderna de ciência.

As fontes consultadas permitem afirmar que o Dresden contém tabelas e registros astronômicos. Milbrath associa ao manuscrito o ciclo de Vênus, de aproximadamente 584 dias, e sua recorrência em um ciclo de oito anos solares (1988, p. 67). Aveni também menciona a tabela de Vênus do códice (2003, p. 162), enquanto Broda registra tabelas de lunações e eclipses (2012, p. 115). Esses testemunhos sustentam a existência de registros de periodicidades celestes. O material reunido para este ensaio, contudo, não detalha suficientemente os procedimentos de correção das tabelas nem permite afirmar que elas previssem a visibilidade de eclipses em um local específico.

É necessário distinguir esse conteúdo do quadro geral da numeração maia. Broda descreve a notação posicional encontrada em inscrições e a Conta Longa (2012, pp. 111, 113); Monteferrante explica os sinais de ponto, barra e concha e a estrutura calendárica que inclui o *tun* de 360 dias (2012). Essas informações ajudam a contextualizar práticas numéricas maias, mas não demonstram, por si mesmas, como cada tabela do Dresden foi calculada ou corrigida. Neste ensaio, portanto, a discussão numérica não será usada como reconstrução do algoritmo do códice.

O formato manuscrito permite examinar como sequências de sinais e imagens organizam informação em uma página. Houston e Taube argumentam, em escala mesoamericana, que escrita e dispositivos gráficos podiam evocar fala e outros sentidos, e que a escrita não era apenas um registro passivo (2000, pp. 261, 265). Essa abordagem oferece uma pergunta interpretativa sobre a legibilidade e a dimensão comunicativa do Dresden; não é uma análise específica de todas as suas páginas e não autoriza atribuir aos escribas uma teoria estética moderna.

O caso documenta, assim, a articulação entre tabelas, periodicidades celestes e registro gráfico, dentro dos limites do manuscrito preservado e das fontes usadas. A ideia de que a organização visual do códice produz uma experiência estética ou torna o cálculo socialmente inteligível é uma interpretação do ensaio. “Estética da Exatidão” nomeia essa lente contemporânea; não constitui uma categoria nativa demonstrada pelas fontes.

\subsection{ Murais de Bonampak: narrativa, percepção e poder}

Os murais de Bonampak são tratados aqui como um programa de pintura narrativa, mas o corpus consultado não inclui uma monografia dedicada à cronologia, à iconografia completa, à técnica pictórica ou à história política do sítio. O plano situa as pinturas no Clássico Tardio; essa datação deve ser sustentada por bibliografia específica antes da versão final, pois os estudos usados abaixo não oferecem, nos trechos auditados, uma demonstração cronológica suficiente. Pela mesma razão, não se fazem afirmações sobre pigmentos, conservação ou sequência construtiva das salas.

Houston e Taube reproduzem uma figura de cantor e o glifo associado na Sala 1, usando o exemplo para discutir relações entre escrita, fala e representação (2000, p. 275, fig. 9; pp. 275–276). O estudo aborda, de forma mais ampla, como dispositivos gráficos mesoamericanos podiam mobilizar sentidos e como a experiência visual era culturalmente organizada (pp. 261, 287–289). Essa análise ajuda a formular perguntas sobre a comunicação dos murais, mas não substitui uma leitura iconográfica específica do programa inteiro de Bonampak. A passagem do exemplo do cantor para afirmações sobre todas as cenas deve, portanto, permanecer explicitamente interpretativa.

Milbrath relata uma hipótese historiográfica que relaciona imagens de guerra e simbolismo de Vênus, incluindo paralelos com Bonampak (1988, pp. 66–67). A autora observa que interpretações astronômicas não anulam necessariamente leituras históricas das cenas (p. 66). Como se trata de uma hipótese secundária apresentada em um estudo voltado à arquitetura e iconografia astronômica de Chichén Itzá, ela não comprova que os murais de Bonampak constituam um calendário ou um registro de observações planetárias. A associação celeste deve ser apresentada como possibilidade debatida e não como resultado estabelecido.

Com as fontes disponíveis, é possível usar Bonampak para discutir a capacidade narrativa e sensorial da imagem e a maneira como representações podem organizar posições e relações sociais. A leitura de legitimação política é plausível, mas requer apoio de pesquisa diretamente dedicada ao sítio para ser desenvolvida como conclusão específica. A comparação com o Dresden também deve ser limitada: o códice documenta tabelas de periodicidades celestes, enquanto a evidência mobilizada para Bonampak trata de imagem, escrita e uma hipótese astral secundária. Os casos permitem comparar suportes e perguntas, mas não demonstram uma equivalência entre suas funções.

\section{ Estética da Exatidão e os limites da categoria moderna de tecnologia}

“Estética da Exatidão” é uma categoria proposta neste ensaio para descrever situações em que organização matemática, repetição calendárica ou orientação espacial participa também de uma experiência visual e coletiva. Ela não significa que a beleza maia derivasse sempre de cálculo, nem que precisão matemática e estética fossem uma unidade conceitual reconhecida pelos próprios produtores. A categoria é útil apenas se mantiver a diferença entre descrição histórica e interpretação contemporânea.

Os três casos mostram modalidades distintas dessa articulação. No Castillo, a forma arquitetônica e a luz solar criam uma experiência pública cujo significado equinocial permanece debatido. No Dresden, cálculo e notação se tornam consultáveis por meio da organização da página. Em Bonampak, pintura, escrita e gesto constroem uma narrativa visual que pode relacionar memória, hierarquia e cosmologia. A precisão, portanto, não é um atributo único compartilhado igualmente pelos três objetos; aparece como possibilidade de ordenamento e comunicação, e deve ser demonstrada em cada suporte.

Esse enquadramento também permite rever o conceito moderno de tecnologia. Se tecnologia for entendida apenas como ferramenta que resolve um problema prático, perde-se a dimensão de representação, transmissão, ritual e autoridade. Se toda prática simbólica for chamada de tecnologia, porém, o termo perde poder explicativo. Neste ensaio, tecnologia designa de modo limitado os meios materiais e as práticas organizadas que permitem observar, registrar, construir e transmitir conhecimento: alinhamentos e edifícios, tabelas e escrita, pintura e sistemas gráficos. Aveni (2003, pp. 174–175) adverte que marcadores temporais podem ser também cenários rituais; Broda (2012, pp. 106–109) situa conhecimento nas instituições e cosmologias; Houston e Taube (2000, pp. 261, 265) mostram a atividade sensorial da escrita e da imagem.

A relação com a transdisciplinaridade contemporânea é uma comparação, não uma continuidade histórica. O trabalho atual separa arqueologia, astronomia, história da arte e antropologia em disciplinas acadêmicas; a reconstrução de práticas antigas exige que essas áreas dialoguem. Mas isso não significa que a sociedade maia possuísse um programa de “transdisciplinaridade” no sentido universitário contemporâneo. O ganho filosófico é reconhecer que nossas divisões analíticas são ferramentas históricas, não fronteiras universais da experiência humana.

\section{ Conclusão}

A revisão sustenta uma versão delimitada da hipótese da Carta de Intenção. Em contextos maias, conhecimentos calendáricos e observacionais foram registrados em manuscritos e relacionados a arquitetura; imagens e espaços participaram de formas de comunicação sensorial, ritual e política. Aveni, Šprajc e Sánchez Nava convergem em que orientação astronômica não se compreende apenas por sua utilidade, mas os últimos mostram também que a evidência de um fenômeno luminoso não determina sozinha sua data ou finalidade. Milbrath oferece interpretações astronômicas e iconográficas importantes para Chichén Itzá, mas várias permanecem hipóteses. Houston e Taube ajudam a interpretar a potência comunicativa de imagens e escrita, sem transformar Bonampak em prova astronômica. Broda fornece a moldura histórica para pensar observação e cosmologia juntas, desde que sua escala mesoamericana não seja confundida com evidência específica de cada sítio maia.

Assim, a articulação entre ciência, arte, religião e política revela menos uma fusão universal do que uma configuração histórica na qual as mesmas práticas e objetos podiam atravessar funções que as disciplinas modernas separam. O fenômeno do Castillo é visível, mas sua interpretação equinocial exata não está assegurada; o Dresden documenta periodicidades e tabelas astronômicas, embora o corpus consultado não permita detalhar sua previsão de visibilidade local; Bonampak evidencia formas visuais e narrativas cuja associação a Vênus permanece conjectural. Essas diferenças de confiança são parte da conclusão, não uma ressalva lateral.

“Estética da Exatidão” contribui como lente contemporânea para examinar como ordem numérica, arquitetura, imagem e experiência pública podem se cruzar. Seu limite é igualmente importante: não deve ser projetada como conceito nativo nem usada para romantizar a sociedade maia. O que esses casos oferecem à reflexão atual é uma crítica às classificações rígidas, acompanhada de um método para não substituí-las por uma idealização igualmente simplificadora: partir dos objetos, identificar as evidências, explicitar as inferências e manter aberta a diferença entre o que se vê e o que se interpreta.

\end{document}
